\chapter{The grid as a physical and numerical model}\label{grid-deep-dive}

A grid does more than select where an equation is evaluated. It decides which
volume owns mass, which surface carries a flux, which neighboring states can
influence one another, and which motions can be represented. These decisions
must agree between dynamics, microphysics, terrain, diagnostics and output.
This chapter develops those agreements from simple cases before considering
stretched grids and nested domains. The derivations below explain constraints
on the implementation; they are not claims that every Nimbus operator is an
exact realization of every identity derived here.

\section{A cell average, a sample and an inventory are different objects}

Let $C$ be a physical cell with volume $V_C$. For a continuous density field,
the cell-average dry density is
\[
 \overline{\rho_d}_C=\frac{1}{V_C}\int_C\rho_d\,dV,
 \qquad M_C=V_C\overline{\rho_d}_C.
\]
If a vapor field varies inside the cell, the dry-mass-weighted mixing ratio is
\[
 \widetilde q_{v,C}
 =\frac{\int_C\rho_d q_v\,dV}{\int_C\rho_d\,dV},
 \qquad W_{v,C}=M_C\widetilde q_{v,C}.
\]
In general, this weighted average is not the arithmetic volume average of
$q_v$, and neither average has to equal the value at the geometric center.
Writing one floating-point number per cell necessarily chooses a representation
of the unresolved variation. Nimbus's conservative scalar transport treats
$\rho_d V$ as the carrier inventory and $Mq$ as the transported species
inventory. A plotted center value is a reconstruction of that representation.

The distinction matters even in an apparently simple two-cell average. Suppose
two equal-volume cells contain dry masses of 2 and 1 kg and have mixing ratios
of 0.01 and 0.02. Their total water is 0.04 kg and total dry mass is 3 kg, so the
combined mixing ratio is $0.04/3$. The arithmetic mean, 0.015, would imply
0.045 kg of water after combining them. No equation of state or more accurate
color map can repair that inventory error afterward. Restriction, remapping
and budget calculations must preserve the mass weights at the operation that
changes the representation.

Nimbus is not uniformly a finite-volume discretization of every prognostic
quantity. Scalar continuity and conservative transport have explicit inventory
contracts; velocity tendencies, pressure projection and closures have their own
discrete forms. The conservation of one scalar update does not establish a
conservative total-energy discretization of the complete model. The useful
question is which identity an operator satisfies, using which state and which
geometry, rather than whether the entire program is given one numerical label.

\section{Allocated dimensions are not the atmospheric domain}

The arrays are contiguous containers with index
\[
 \ell(i,j,k)=(iN_y+j)N_z+k.
\]
The last index is contiguous in memory. This is a storage rule, not a definition
of a cell's physical volume or a promise that each index is an independent
atmospheric cell. The historical member names \texttt{nr} and \texttt{nth}
are also used by Cartesian paths; their spelling does not make those paths
cylindrical. Coordinates, staggering and the declared physical domain supply
the interpretation.

\href{../include/core/infra/physical_cell_domain.hpp}{PhysicalCellDomain}
uses half-open ranges, so $i_{\rm begin}\le i<i_{\rm end}$ owns scalar
control volumes. In the ordinary arrays with at least three points on an axis,
one horizontal plane at each side is boundary-closure storage. The collocated
layout similarly excludes one plane at each vertical end when there are at
least three vertical entries. The C-grid scalar layout retains its vertical
scalar layers, with its ground-face velocity handled separately. Tiny arrays
used by bounded tests have explicit special cases; do not derive a universal
halo width from a single fixture.

For example, an allocated $10\times8\times6$ scalar array has 480 storage
locations. Under the C-grid ownership rule it contains $8\times6\times6=288$
physical scalar cells. Under the ordinary collocated rule it contains
$8\times6\times4=192$. These counts describe the current ownership helper,
not recommended simulation dimensions. Multiplying 480 by a nominal cell
volume would count closure state as additional atmosphere in either example.
Likewise, inferring the physical domain length from the allocation alone can
change a prescribed wavelength, mean density or column depth.

A halo can contain a finite, plausible value copied from an interior cell.
That does not give it volume. It exists so a stencil can read a defined neighbor,
so a boundary algorithm can retain history, or so a restart can reproduce the
same closure state. A valid restart may therefore retain more data than a
physical inventory integrates. Output metadata must carry the ownership range
rather than asking a reader to guess it from zeros, NaNs or array dimensions.

\section{Three distances that must not be collapsed into one}

Consider a flat but vertically stretched column. Let the lower and upper
interfaces of layer $k$ be $z_{k-1/2}$ and $z_{k+1/2}$, and let its representative
scalar height be $z_k$. The layer thickness is
\[
 h_k=z_{k+1/2}-z_{k-1/2},
\]
whereas the distance used by a center-to-center pressure gradient is
$d_{k+1/2}=z_{k+1}-z_k$. A vertical-velocity dual control volume spanning two
adjacent scalar centers has a width related to $d_{k+1/2}$, not automatically
to either neighboring $h_k$. All three distances coincide in a special uniform
construction; that coincidence is the reason the distinction can disappear
unnoticed in a uniform-grid implementation.

For a piecewise linear reconstruction between two scalar locations, the value
at an intermediate interface $z_f$ is
\[
 \phi_f=\frac{z_{k+1}-z_f}{z_{k+1}-z_k}\phi_k
       +\frac{z_f-z_k}{z_{k+1}-z_k}\phi_{k+1}.
\]
A simple half-sum is justified only when the face is halfway between those
locations. An independently reconstructed upwind value is a different numerical
choice. Neither should be substituted for the other just because both functions
return a value at an interface.

The divergence of a vertical flux density $f$ in layer $k$ is
$(f_{k+1/2}-f_{k-1/2})/h_k$. The gradient at an interior interface instead uses
$(\phi_{k+1}-\phi_k)/d_{k+1/2}$. Substituting one global $\Delta z$ into both
expressions changes the amount of transport per layer and the force between
layers. It can also destroy the cancellation between a base pressure gradient
and gravity while leaving the fields visually smooth.

The \href{../include/core/infra/physical_grid_metrics.hpp}{physical metric
interface} validates scalar centers, interfaces and positive layer thicknesses
and supplies the geometry used by column physics. It rejects a partly initialized
flat geometry and a terrain geometry whose shape does not match the fields.
Radiation, sedimentation, vertical mixing and diagnostics must consume that
same physical column. A radiation heating rate based on a nominal layer depth
while transport uses the actual depth is a coupling inconsistency, not merely
a small plotting discrepancy.

\section{Why C-grid staggering changes resolved waves}

The C-grid locates a normal velocity at the face through which it transports
mass. In Nimbus's index convention, $u(i,j,k)$ is the east face of scalar cell
$(i,j,k)$, $v(i,j,k)$ its north face and $w(i,j,k)$ its upper face. Thus the west
face used by a cell is addressed with the neighboring index. The separately
owned lower ground face is an important endpoint of the vertical construction.
A physical scalar interval of $n$ cells has $n+1$ boundary-and-interior face
locations in a nonperiodic direction. Under periodic conditions the end faces
represent one physical connection and must not become two distinct exchanges
in an inventory sum.

The benefit can be seen without a storm simulation. For a uniform one-dimensional
pressure perturbation $p_j=\widehat p\exp(\mathrm{i}kj\Delta x)$, a collocated
centered gradient has the Fourier factor
\[
 \frac{p_{j+1}-p_{j-1}}{2\Delta x}
 =\mathrm{i}\frac{\sin(k\Delta x)}{\Delta x}p_j.
\]
It is blind to the alternating grid-scale pattern at $k\Delta x=\pi$. A
center-to-face C-grid difference instead has magnitude
$2\sin(k\Delta x/2)/\Delta x$ after accounting for the half-cell phase shift.
It does not vanish at that wavelength. This explains one pressure--velocity
coupling advantage of staggering. It does not imply that a two-cell wave is
accurate or that an A-grid implementation cannot provide a different consistent
coupling with its own stabilization and reconstruction.

For the illustrative linear acoustic pair
$p_t=-\rho_0c^2 D u$ and $u_t=-G p/\rho_0$, a compatible C-grid divergence
and gradient on a periodic uniform grid give
\[
 \omega_h^2=c^2\left[\frac{2\sin(k\Delta x/2)}{\Delta x}\right]^2.
\]
The continuous relation is $\omega^2=c^2k^2$. Their difference is a spatial
dispersion error before the time integrator contributes any additional error.
At small $k\Delta x$, expanding the sine shows convergence to the continuous
wave; near the shortest represented wavelength, the discrete phase and group
velocities differ substantially. Increasing the formal order of an unrelated
microphysics source does not repair that wave dispersion.

A nominal grid spacing therefore does not equal the smallest accurately
resolved physical structure. The useful resolved range depends on the complete
spatial operator, filtering, time step, closures and the diagnostic being
measured. A resolution study must compare that diagnostic across grids under
consistent physical assumptions. Holding a number of grid cells fixed while
changing $\Delta x$ changes the physical size of a perturbation; holding a
Smagorinsky coefficient fixed while changing its grid-derived filter width
changes the modeled subgrid diffusivity.

\section{Flux cancellation is a topology statement}

Assign each interior face one orientation and one numerical dry-mass flow $F_f$.
Let $B_{Cf}$ be $+1$ when that orientation points outward from cell $C$, $-1$
when it points inward, and zero otherwise. The update is
\[
 M_C^{n+1}=M_C^n-\Delta t\sum_f B_{Cf}F_f.
\]
An interior column of $B$ contains one $+1$ and one $-1$. Summing over all
physical cells cancels that face exactly in real arithmetic, regardless of
whether the reconstruction is first or fifth order. What remains is the
physical boundary flux. This is the fundamental topological reason shared
fluxes conserve an inventory.

If the two adjacent cells separately reconstruct and round what is supposed to
be the same interface flow, they can use $F_f+\epsilon_L$ and
$F_f+\epsilon_R$. The domain acquires a residual proportional to
$\epsilon_L-\epsilon_R$. Thread safety alone does not establish a shared flux:
two race-free calculations can still disagree. Conversely, deterministic
summing of a wrong set of volumes does not make an inventory physical.

For a mixing ratio, Nimbus's retained continuity flux provides the same carrier
flow to the species update,
\[
 Q_C^{n+1}=Q_C^n-\Delta t\sum_f B_{Cf}F_f q_f,
 \qquad q_C^{n+1}=Q_C^{n+1}/M_C^{n+1}.
\]
When $q_f=q_0$ everywhere, the uniform value survives exactly in real arithmetic:
$Q_C^{n+1}=q_0 M_C^{n+1}$. This is stronger than obtaining the correct domain
sum alone. A scheme can preserve total water yet generate false cell-scale
concentration changes if it transports a different dry-mass carrier. The
retained-flux transport contracts test this consistency along with positivity
and invalid-input behavior.

The algebra also exposes a time-level requirement. If $Q_C^n$ is constructed
from an already updated density and an old mixing ratio, it is not the old
inventory. Saving the old carrier, or constructing the inventory before the
carrier is overwritten, is part of the mathematical scheme. Eliminating that
snapshot for performance requires an equivalent ownership arrangement. It is
not safe merely because no API signature changes.

\section{Reconstruction determines accuracy; limiting determines admissibility}

A face flow needs a transported value. In a positive one-dimensional flow, a
first-order upwind method takes the value from the cell on the inflow side.
For constant velocity $a>0$ and uniform spacing, the update is
\[
 q_j^{n+1}=(1-C)q_j^n+Cq_{j-1}^n,
 \qquad C=a\Delta t/\Delta x.
\]
For $0\le C\le1$, the new value is a convex combination and cannot exceed the
two contributing extrema. This explanatory result assumes a constant carrier,
no source and the stated first-order discretization. It is not a universal
positivity proof for a compressible, multidimensional, source-coupled update.

A linear reconstruction uses a slope to represent subcell variation and
reconstruct left and right interface states. An unlimited centered slope is
accurate on a smooth field but may produce negative condensate next to a sharp
edge. A limiter reduces the slope when the neighboring differences suggest an
extremum or discontinuity. WENO methods instead combine candidate
reconstructions with smoothness-dependent weights. Their motivation is to
recover high order in smooth regions while avoiding an oscillatory stencil
near a sharp feature. The reconstruction, its limiting, its boundary stencil
and its time integration must be evaluated together.

Nimbus's \href{../src/numerics}{numerical implementations} and the retained
\href{../tests/numerics}{numerical tests} define the admitted transport paths.
The method name does not establish high-order convergence everywhere: a limiter
can intentionally revert toward lower order, a boundary may use a shorter
stencil, and a phase source can introduce nonsmooth behavior. A resolution
comparison should report where the observed error comes from, rather than
assigning the formal smooth-field order to an entire convective storm.

Clipping a negative mixing ratio after a flux update changes its inventory.
If a stored cell would contain $Mq=-\delta$ and is simply set to zero, the
operation adds $\delta$ to the domain water amount. Small clipping is therefore
not automatically harmless. A conservative correction must identify where the
corresponding amount is removed, respect the available donor inventory, and
report the adjustment. A physical source limiter has a related duty: multiple
processes cannot each consume the same full reservoir independently.

\section{Pressure work needs compatible scalar and velocity geometry}

On a uniform grid with periodic boundaries, summation by parts gives a discrete
counterpart of the continuous pressure-work identity. For an interior face
between left and right cells with area $A_f$ and center distance $d_f$,
\[
 \sum_C V_C p_C(Du)_C
 =\sum_f A_f u_f(p_L-p_R)
 =-\sum_f A_f d_f u_f(Gp)_f.
\]
The velocity dual volume is $V_f^*=A_fd_f$. At a physical boundary there is an
additional boundary pressure-work term. This identity depends on using the same
face orientation, areas and center distances on both sides. Interpolating a
pressure gradient with one geometry and diagnosing work with another breaks
the pairing.

If the velocity equation uses $\rho_f\,du_f/dt=-(Gp)_f$, multiplying by
$V_f^*u_f$ supplies the kinetic-energy pressure work. Its partner in the internal
energy equation has the opposite sign. This is a design criterion for a coupled
energy discretization. Nimbus's diagnosed energy budgets and pressure/EOS
operators must be assessed against their actual discrete contracts; the
continuous identity by itself does not prove exact total-energy conservation
for the implemented rho--theta model.

A staggered kinetic inventory must first decide which cell mass weights its
velocity locations. As an illustrative dry inventory, retain the dry mass $M_C$
defined above. On a flat grid, assigning half of each adjacent cell's mass gives
\[
 M_f^*=\tfrac12M_L+\tfrac12M_R,
 \qquad K_f=\tfrac12M_f^*u_f^2.
\]
Boundary faces, vertical stretching and mapped ground motion require the
appropriate half-volume and metric construction. Averaging velocity to scalar
centers and then squaring is a different functional because the square of an
average is not the average of squares. Nimbus's moist-mixture basis-v2 diagnostic
instead uses cell masses $M_C^{\rm mix}=\rho_{d,C}V_C(1+q_{t,C})$ in this
half-cell construction. That diagnostic choice does not change the dynamical
pressure-force convention, which uses its declared dry-density approximation.
The flat-grid
\href{../src/diagnostics/budgets/velocity_energy_pair.cpp}{velocity/energy
pairing implementation} and the full
\href{../src/diagnostics/budgets/energy_budget.cpp}{energy diagnostic} make the
mass basis and geometric scope explicit.

Time discretization adds another issue. At fixed face mass, an impulse $I$
changes velocity by $I/M_f^*$ and changes kinetic energy by
\[
 \Delta K_f=u_f^n I+\frac{I^2}{2M_f^*}.
\]
The second term is missing if work is calculated using only the old velocity.
When face mass also changes, a term for that mass change is required. For
example, the exact endpoint decomposition
\[
 \Delta K_f=\tfrac12 M_f^{*,n+1}
  \big[(u_f^{n+1})^2-(u_f^n)^2\big]
 +\tfrac12\big(M_f^{*,n+1}-M_f^{*,n}\big)(u_f^n)^2
\]
is algebraically valid but assigns work differently from a midpoint-based
construction. An energy report must declare its endpoint basis; combining
terms from two different decompositions can double-count a contribution.

\section{Terrain is a coordinate transformation, not a tilted picture}

For the explanatory mapping $(x,y,\zeta)\mapsto(x,y,z(x,y,\zeta))$, define
$J=z_\zeta>0$, $z_x=(\partial z/\partial x)_\zeta$ and
$z_y=(\partial z/\partial y)_\zeta$. Chain rules give
\[
 \begin{aligned}
 (\partial_x\phi)_z&=(\partial_x\phi)_\zeta
                         -\frac{z_x}{J}\partial_\zeta\phi,\\
 (\partial_y\phi)_z&=(\partial_y\phi)_\zeta
                         -\frac{z_y}{J}\partial_\zeta\phi,\\
 \partial_z\phi&=\frac{1}{J}\partial_\zeta\phi.
 \end{aligned}
\]
Nimbus's mapped transport interface uses $G_x=-z_x$ and $G_y=-z_y$ with this
Jacobian convention. The signs are part of the interface. Reusing a formula
from a paper whose metric symbols have the opposite sign without converting
it can reverse the terrain correction while preserving dimensions.

For a stationary mapping the contravariant vertical coordinate velocity is
\[
 \dot\zeta=\frac{w-u z_x-v z_y}{J}.
\]
The transformed dry-mass equation can be written
\[
 \partial_t(J\rho_d)+\partial_x(J\rho_d u)
 +\partial_y(J\rho_d v)
 +\partial_\zeta\{\rho_d(w-u z_x-v z_y)\}=0.
\]
This equation is written for the stated mapping and continuous metrics. A
finite-volume implementation must still construct compatible discrete face
metrics, fluxes and physical volumes. It cannot obtain that compatibility merely
by evaluating the continuous expressions independently at arbitrary points.

At stationary terrain $z=h(x,y)$, zero normal flow requires
$w=u h_x+v h_y$. Setting the Cartesian vertical velocity to zero instead imposes
a different condition unless the terrain is locally flat or the wind has no
component along the slope gradient. A column's lower interface can therefore
have a physical vertical-velocity contribution even when no mass crosses the
ground. The distinction carries into surface stress, terrain pressure work and
the staggered kinetic inventory.

A useful geometric consistency check is free-stream preservation. A constant
physical velocity field should not acquire spurious divergence solely because
coordinate surfaces slope. The discrete metric construction must allow the
corresponding face contributions to cancel. A second check is hydrostatic
balance: pressure and gravity should cancel for the declared resting base
state using the same physical heights and averaging. These are different tests.
Passing one does not prove the other or establish the accuracy of transport
through a sharp terrain transition.

See the \href{../include/terrain/conservative_transport.hpp}{mapped gradient
and transport contract},
\href{../src/terrain/base/conservative_transport.cpp}{its implementation}, and
\href{../tests/dynamics/test_mapped_reconstruction.cpp}{mapped reconstruction
checks}. The \href{https://www2.mmm.ucar.edu/wrf/site/documentation/technote.html}{WRF
technical notes} provide a useful comparison with a mature terrain-following
solver. WRF's vertical coordinate and complete equation set are not Nimbus's;
that reference motivates questions about metrics and staggering rather than
supplying an interchangeable formula set.

\section{A moving window is not a new physical wind}

For a moving physical control volume, Reynolds transport uses relative normal
velocity $(\boldsymbol u-\boldsymbol v_g)\cdot\boldsymbol n$. If the window
translates horizontally with velocity $\boldsymbol U_g$, a stationary Earth-frame
parcel flows backward through that window even though its physical wind is
zero. Stored Earth-frame momentum must not be incremented by a window shift.
The relative transport velocity and the physical velocity used for kinetic
energy have different purposes.

A continuously deforming coordinate system also obeys the geometric conservation
law: the change in its volume must match the surface integral of grid velocity.
That is a constraint on a general moving-grid discretization, not a claim that
Nimbus implements unrestricted mesh deformation. Nimbus's domain controls,
integer relocation, terrain placement and supported nesting profiles have
specific execution and admission contracts. A discrete relocation is a remap
operation with retained and newly exposed regions; it is not equivalent to
integrating a continuous ALE flux over an arbitrary deformation.

To understand a conservative remap, let $m_{ij}$ be dry mass transferred from old
cell $j$ into new cell $i$. The transferred species amount is $m_{ij}q_j$ in a
piecewise-constant example. Then
\[
 M_i^{\rm new}=\sum_jm_{ij},\qquad
 Q_i^{\rm new}=\sum_jm_{ij}q_j,\qquad
 q_i^{\rm new}=Q_i^{\rm new}/M_i^{\rm new}.
\]
If the transfers are nonnegative, the reconstructed ratio is bounded by its
donors. If $\sum_i m_{ij}=M_j^{\rm old}$ for every fully retained old cell,
the remap conserves its mass. Cells leaving a limited window and fresh cells
entering it require explicit departure and inflow accounting. The mass inside
a tracking window need not remain constant even when the remap is correct.

Private particle moments, held boundary-layer sources, turbulence state and
boundary history also have location and time ownership. Preserving the visible
$q$ arrays while dropping a number moment changes the size distribution.
Preserving a held heating array while assigning it to the wrong Earth location
changes the forcing. The
\href{../src/core/infra/domain_recycling.cpp}{domain-recycling implementation}
and its provenance records identify the geometric ownership that a relocation
must preserve.

\section{Nesting introduces an interface budget}

A fine child overlaps part of a coarse parent. Adding every parent cell and every
child cell counts the overlap twice. A composite inventory instead includes
the uncovered parent plus the child representation of the covered region.
This choice must be explicit before asking whether two-way coupling conserves
mass or energy. A domain total from one grid and a composite total are different
diagnostics even if they use the same units.

For restriction of a conserved scalar, the coarse covered inventory is the sum
of its contributing child inventories. The coarse mixing ratio follows by
dividing summed species mass by summed dry mass. Face-normal momentum and
staggered velocity require their own geometry and cannot be restricted by
copying the scalar formula blindly. Prolongation must likewise distinguish
preserving a parent mean from creating artificial fine-grid extrema.

Flux mismatch is a separate concern. Suppose the coarse update removes
$\Delta tF_c$ at a coarse/fine interface, while fine substeps remove a total
$\sum_s\Delta t_s\sum_{f\subset c}F_{f,s}$. A correction replacing the coarse
exchange with the fine exchange has the illustrative signed form
\[
 \delta M_c=\Delta t F_c-
            \sum_s\Delta t_s\sum_{f\subset c}F_{f,s}.
\]
The sign reverses with the orientation convention. Matching instantaneous
fluxes at the final time is not enough when the child subcycles. Water and
thermal carriers also need consistent exchanges; matching dry mass alone does
not establish the coupled conservation result.

Nimbus separates nested level execution, scalar feedback, terrain geometry and
boundary ownership. The
\href{../include/core/infra/nested_scalar_feedback.hpp}{feedback contract},
\href{../src/core/infra/nested_scalar_feedback.cpp}{feedback implementation} and
\href{../src/core/infra/nested_execution.cpp}{level execution} are the relevant
source boundaries. The interface arguments above explain what must be checked;
they do not assert that every nested configuration implements this particular
reflux formula or that arbitrary combinations are admitted. The public runtime
catalog remains authoritative about the supported combination.

\section{Spatial resolution also changes the time and closure problem}

An advective time scale in a cell is set by its outgoing relative flux compared
with its carrier inventory. A useful dimensional measure is
\[
 C_M=\frac{\Delta t}{M_C}\sum_{f\,\mathrm{out}}|F_f|.
\]
For simple upwind transport this is directly related to positivity. The exact
stability bound depends on the spatial operator and time integrator. Acoustic
steps have the additional sound-speed scale; explicit diffusion has a scale
proportional to $h^2/K$; sedimentation has a fall-speed crossing time. Improving
one constraint does not remove the others.

On a uniform Cartesian mesh, familiar explanatory estimates are
\[
 \Delta t_{\rm adv}\sim\frac{h}{U},\qquad
 \Delta t_{\rm acoustic}\sim\frac{h}{c+U},\qquad
 \Delta t_{\rm diff}\sim\frac{h^2}{K}.
\]
The constants and directional combinations must come from the implemented
scheme, not these dimensional estimates. Mapped grids must use the relevant
physical metric; a single nominal horizontal spacing does not detect a very
thin vertical layer. Nimbus exposes acoustic substep selection and time-step
history through the \href{../include/numerics/time_stepping/time_stepping_base.hpp}{time-stepping
contract}, with separate stability and cadence owners elsewhere in the runtime.

Refining all three spatial axes by two creates roughly eight times as many
physical cells. If a wave or advective constraint also halves the time step,
work per simulated interval grows by about sixteen before changes in solver
iterations, output, memory traffic or physics subcycling. This is a scaling
argument, not a benchmark prediction. If explicit diffusion becomes limiting,
the time-step penalty can be larger. Refinement also changes grid-derived SGS
filter widths and the relation between resolved clouds and bulk closures.
A scientifically useful convergence experiment must specify which physical
parameters stay fixed and which model assumptions change with resolution.
